\section{Introduction}\label{main:introduction}

Bose--Einstein condensation means that a single one-particle state carries
a macroscopic part of a many-particle system. For a homogeneous gas the
natural candidate is the constant orbital. We prove that this orbital
contains a fixed positive fraction of every ground state of a sufficiently
dilute three-dimensional hard-sphere Bose gas, in the thermodynamic limit
at fixed density and fixed exclusion distance. The conclusion permits
complex wave functions and a degenerate ground space.

\subsection{The result}

Let $\Lambda_L=(\R/L\Z)^3$ be the three-dimensional flat torus, with
its Euclidean torus distance $d_L$. For $N$ particles and exclusion
distance $a>0$, define the open configuration set
\[
 \Omega_{N,L,a}
  =\{X\in\Lambda_L^N:d_L(x_i,x_j)>a\text{ whenever }i<j\}.
\]
The hard-sphere Hamiltonian is the operator associated with the Dirichlet
form
\[
 q_{N,L,a}[f]=\sum_{i=1}^N\int_{\Omega_{N,L,a}}|\nabla_i f|^2,
 \qquad f\in H_0^1(\Omega_{N,L,a}),
\]
restricted to symmetric functions. Thus its kinetic energy is
$-\sum_i\Delta_i$, and the wave function vanishes at forbidden
configurations. The exclusion distance $a$ is the scattering length:
the exterior zero-energy scattering solution is $1-a/|x|$.
All wave functions below are extended by zero to $\Lambda_L^N$.

For a normalized bosonic vector $\Psi$, its one-particle density matrix
and the normalized constant orbital are
\[
 \Gamma^{(1)}_\Psi(x;y)
   =N\int_{\Lambda_L^{N-1}}
           \Psi(x,Y)\overline{\Psi(y,Y)}\,\dd Y,
 \qquad \varphi_{0,L}=L^{-3/2}.
\]
In particular, $\Tr\Gamma^{(1)}_\Psi=N$. We measure condensation by
\begin{equation}\label{main:occupation}
 B(\Psi)=\frac{\langle\varphi_{0,L},
                      \Gamma^{(1)}_\Psi\varphi_{0,L}\rangle}{N}
        =\frac1{L^3}\int\left|\int\Psi(x,Y)\,\dd x\right|^2\dd Y.
\end{equation}
This quantity lies in $[0,1]$.

\begin{theorem}[Condensation in every ground state]\label{main:theorem}
There are absolute constants $\varepsilon_0,c_0>0$ with the following
property. Fix $a,\rho>0$ satisfying $\rho a^3<\varepsilon_0$.
For every sequence of particle numbers $N_k$ and lengths $L_k$ such that
\[
 N_k\longrightarrow\infty,\qquad L_k\longrightarrow\infty,
 \qquad \frac{N_k}{L_k^3}\longrightarrow\rho,
\]
and every choice of normalized, possibly complex, bosonic ground
eigenvectors $\Psi_k$ of the hard-sphere Hamiltonians,
\begin{equation}\label{main:limit}
 \liminf_{k\to\infty} B(\Psi_k)\ge c_0.
\end{equation}
\end{theorem}

The same constant $c_0$ works at all sufficiently small fixed gas
parameters $\rho a^3$. The volume needed to approach this bound may
depend on $a$ and $\rho$. We prove a positive fraction; we do not
estimate the optimal depletion or show that the fraction tends to one
as $\rho a^3$ tends to zero.

\begin{corollary}[Ground-space states]\label{main:mixed}
Under the hypotheses of Theorem~\ref{main:theorem}, the same lower
bound holds for every sequence of density operators supported on the
bosonic ground eigenspaces. In particular, it holds for the normalized
ground-space projections $\Pi_k/\Tr\Pi_k$, and for every sequence of
normalized nonnegative ground functions.
\end{corollary}

Here the one-particle density matrix of a density operator is defined by
linearity. The proof below is uniform over its pure ground vectors;
convexity then proves the corollary. No choice of a preferred ground
vector or a basis of positive ground functions is required.

\subsection{Context}

The ideal-gas prediction grew from Bose's counting of light quanta and
Einstein's extension to a gas of material particles
\cite{Bose1924,Einstein1924IdealGas,Einstein1925IdealGas}.
For an interacting gas, the density-matrix criterion of Penrose and
Onsager~\cite{PenroseOnsager1956} identifies condensation through a
macroscopic one-particle occupation. Bogoliubov's theory assumes a
macroscopically occupied zero mode to derive the weakly interacting gas's
excitation picture~\cite[pp.~24--25]{Bogoliubov1947}. For hard spheres,
Lee, Huang and Yang's pseudopotential calculation predicts a depleted
fraction of order $\sqrt{\rho a^3}$
\cite[Section~III]{LeeHuangYang1957}. The rigorous occupation problem is
to justify macroscopic condensation itself before seeking this sharper
depletion law.

The hard-sphere energy theory runs from Dyson's bounds~\cite{Dyson1957}
and the leading dilute asymptotics of Lieb and Yngvason
\cite{LiebYngvason1998} to the sharp Lee--Huang--Yang lower bound of
Fournais and Solovej, which includes hard cores
\cite{FournaisSolovej2023}. Basti, Brooks, Cenatiempo, Olgiati and
Schlein proved the matching upper coefficient for hard spheres and
subsequently for a class of general nonnegative potentials
\cite{BastiEtAl2026HardSpheres,BastiEtAl2026GeneralPotentials}.
These energy asymptotics do not determine single-orbital occupation as
the one-particle spectral gap vanishes in large volume. Lieb and
Yngvason explicitly separated that question from their energy theorem
\cite{LiebYngvason1998}.

Continuum condensation results allowing scaled hard cores include the
trapped Gross--Pitaevskii theorem of Lieb and Seiringer
\cite{LiebSeiringer2002} and the homogeneous canonical
positive-temperature theorem of Deuchert and Seiringer
\cite{DeuchertSeiringer2020GP}. Low-energy results of Fournais and of
Chong, Liang and Nam~\cite{Fournais2021,ChongLiangNam2026}, and Junge's
canonical low-temperature result in Neumann boxes
\cite[Corollary~6]{Junge2026}, treat joint dilute and increasing-volume
limits. These regimes permit the gas parameter to decrease as the volume
grows. Here $a$ and a positive $\rho$ are fixed before the volume grows
without restriction, and the lower fraction is uniform over all ground
vectors.

S\"ut\H{o}'s thermodynamic condensation theorem, above a temperature-dependent
density threshold, assumes nonnegative potentials with nonnegative Fourier
transform and additional regularity and integrability
\cite[Theorem~1.2]{Suto2023TwoStep}; these hypotheses exclude a literal hard
core. Galanda and Pinamonti construct algebraic equilibrium states for smooth
interactions, using coupled cutoff-removal and weak-interaction limits in
one setting and Gross--Pitaevskii scaling in another
\cite{GalandaPinamonti2025Equilibrium,GalandaPinamonti2026GP}. Their state
formulations and limit orders differ from the present hard-sphere problem.
Solovej's 2025 survey identifies fixed-density ground-state condensation
as a major open problem~\cite[Section~5]{Solovej2025}.
Theorem~\ref{main:theorem} establishes a positive lower limiting occupation
fraction for sufficiently dilute three-dimensional hard spheres, with the
fixed-parameter order of limits in Equation~\eqref{main:limit}.

The path method has separate antecedents. Benjamini, Pemantle and Peres
constructed unpredictable paths with controlled intersections
\cite[Theorem~1.3 and Lemma~3.1]{BPP1998}. Abbe, Massouli\'e, Montanari, Sly and
Srivastava averaged noisy group differences along paths joining prescribed
endpoints~\cite{AbbeMassoulieMontanariSlySrivastava2018}; Garban and Spencer
adapted this averaging method to symmetry breaking in classical disordered
spin systems~\cite{GarbanSpencer2023}. Averaging path likelihoods before
taking their second moments also underlies trail detection
\cite{AriasCastroCandesHelgasonZeitouni2008,BergerPeres2013} and the
point-reassignment argument of Peres and Sly
\cite[Proposition~2.1]{PeresSly2014}. Our exterior-shell detours use Kesten's
boundary-connectivity theorem in Tim\'ar's formulation
\cite[arXiv version~2, Theorem~4]{Timar2013}.

Our sole companion input is the explicit random lattice skeleton and both
encounter moments in \emph{A density-uniform condensate bound for dilute
Bose gases}~\cite[Section~8.1, pp.~22--24; Lemma~8.1 and
Equations~(24)--(25), p.~23]{SoftCompanion}. These statements contain no
physical parameters. The hard-core energy bounds, obstacle adaptation,
conditional-normalizer calculation and phase argument are proved here.

\subsection{How the proof works}

Write $\Phi=|\Psi|$. The modulus is again a ground function, and we
first seek a lower bound for $B(\Phi)$. Expanding the occupation formula gives
\[
 B(\Phi)=\frac1{L^3}\int\Phi(x,Y)\Phi(y,Y)\,\dd x\,\dd y\,\dd Y.
\]
The same bath $Y$ occurs at the two tagged positions $x$ and $y$. After
partitioning the box into cells, we will obtain a uniformly positive
contribution to this formula from a fixed
positive fraction of ordered cell pairs whose separation is comparable to
the box size. A separate phase estimate will then recover $B(\Psi)$.

Local trajectory changes provide the comparisons between neighboring
cells. Multiplying their costs along a long route would lose a factor at
every step. Instead we average changes over random routes and bound the
second moment of the averaged likelihood. Exact cancellation will leave a
cost determined by mutual encounters of two routes, with a bound independent
of their lengths.

\paragraph{Local trajectories and the exposed environment.}
Change length units so that each unit cell contains on average $D$
particles and the exclusion distance is $r=\alpha/D$, with $\alpha$ in
a fixed interval. Dilution then corresponds to large $D$;
Equation~\eqref{setup:gas-parameter} records the invariant gas parameter.
Local kinetic inequalities and normalized particle insertion show that
cells with too few particles or too many close neighbors have small
\emph{joint} probabilities. An exact stationary killed-Brownian law,
whose position marginal is $\Phi^2$, turns these estimates into a supply
of short trajectories that can be replaced through a home cell or a
face-neighbor cell.

At each cell we allocate the movable labels to many disjoint groups.
For one group at a time, expose the outer endpoints and all trajectories
outside the group. A cell is usable when enough of its remaining labels
admit successful replacements against these frozen trajectories. The
unexposed law is a product of bridge laws that already impose individual
oscillation and confinement cutoffs and avoidance of the frozen paths,
conditioned further on mutual hard avoidance. This exact conditional law, and an unconditional joint
bound on unusable cells, are the two outputs of
Section~\ref{group:section}. We keep the two probability statements
separate throughout the comparison.

\paragraph{Distant cells and a common bath.}
For each required pair of distant cells, the joint obstacle bounds retain
a set of exposed environments of fixed positive probability. On this set,
draw a random lattice route from the companion's skeleton law and reroute
it around encountered finite unusable clusters along their good exterior
shells. The resulting route uses nearest-neighbor steps, so every step
admits the local replacements. Averaging the retained distribution of environments, with two independent
route choices drawn in each common environment, gives a bounded exponential
moment for their mutual encounters. One imported estimate controls encounters of the nominal skeletons; the other
controls the additional encounters caused by detours, after the local
obstacle probabilities have been summed.

Along a route, choose one group label from each departure cell. Construct
two copies of the selected trajectories, preserving each label's two outer
endpoints. Its midpoint lies in its departure cell in the first copy and
in the next cell in the second. At every intermediate cell the two different
labels share one midpoint. Removing the initial tag from the first copy
and the terminal tag from the second leaves the same bath of $N-1$
midpoints; Figure~\ref{cmp:shift-figure} depicts these assignments.
This produces a common finite measure for two position measures that place
a tagged particle in the respective endpoint cells and retain the same
observed bath. Their square-root overlap contributes to $B(\Phi)$ through the exact cell-average identity in
Equation~\eqref{cmp:cell-affinity}.

To lower-bound this overlap, average the change of measure over routes
and label choices, and bound its second moment. The interaction
normalizers are retained exactly. In the product of two likelihoods,
the normalizers from portions of the routes far from one another cancel;
only nearby encounters and shared labels contribute a cost. The preceding
encounter moment controls that cost uniformly in the endpoint separation.
Section~\ref{cmp:section} constructs the two-copy measure before proving
and applying this cancellation. Finally, summing the overlap contributions
of a number of disjoint groups proportional to $D$ compensates for the
normalization of one tagged particle and yields a fixed lower bound for
$B(\Phi)$.

\paragraph{Recovering the phases.}
Equality of the kinetic energies of $\Psi$ and $\Phi$ forces a constant
phase on each connected component of the allowed configuration set.
Those phases may differ between components. Fix the bath and cover its
forbidden balls by small closed cubes, filling the bounded holes of each
connected cluster. The points left outside the filled regions communicate
through the open allowed one-particle slice and hence share one phase.
Small clusters have little total filled volume. A large cluster would
force many nearby particles with close partners, an event controlled by
the joint spatial estimates. If $\delta_{\rm geom}$ is the discarded
volume fraction averaged under the ground-state bath marginal, the
elementary estimate
\[
 B(\Psi)\ge B(\Phi)-4\sqrt{\delta_{\rm geom}}
\]
transfers the modulus bound to the original vector. We first fix the
positive modulus fraction and only then strengthen the dilution threshold
to make this phase loss smaller.

Section~\ref{setup:section} gives the form and phase facts.
Sections~\ref{sp:section}--\ref{bridge:section} supply the spatial, path
and bridge estimates for the local replacements.
Section~\ref{group:section} organizes the groups and their conditional law,
Section~\ref{lat:section} constructs the routes, and
Section~\ref{cmp:section} proves the modulus bound.
Section~\ref{geom:section} controls the filled regions.
Section~\ref{main:conclusion} fixes the constants and returns to arbitrary
thermodynamic sequences in the original units.
